\section{Additional Applications}
\label{app:additional-applications}

The examples below show how the same two questions recur in classical
distributed tasks: can each choice survive more input, and can independent
choices survive together?  We isolate those obstructions here.  The
results do not replace model-specific theorems about failures, computability, or
termination.

\subsection{One commitment, then monotone authority}

Agreement protocols illustrate one distributed form of coordination.  Before an
authority is established, several incompatible authority lineages may remain
possible.  Publishing one lineage rules the others out: that is a commitment.
Suppose the commitment fixes both the initial authority and a rule that uniquely
determines every certified successor.  The system can then expose three kinds
of growing evidence: a certified epoch chain $E$, votes or certificates $V$,
and decisions $D$ indexed by epoch and slot.  Define refinement by
        \begin{equation}
  (E,V,D)\OutputOrd(E',V',D')
  \quad\Longleftrightarrow\quad
  E\preceq_{\rm pref}E',\quad V\subseteq V',\quad D\subseteq D'.
  \label{eq:authority-refinement}
        \end{equation}
Protocol safety ensures that a later certificate extends the unique lineage and
preserves prior decisions.  Public authority changes are therefore refinements,
not revisions: votes, certificates, epochs, and decisions all accumulate after
the initial commitment.

This statement does not claim that later certificates are free.  Producing one
may require communication and agreement.  It says something more specific:
later protocol work does not choose among remaining output alternatives or
retract a public result.  If the first commitment did not fix a unique
succession rule, competing next epochs would present another output choice
and require another commitment.

\subsection{Why $k$-set agreement needs finite families}

Consensus forbids two different decision values.  The more general $k$-set
agreement task allows at most $k$ values, which makes pairwise examples
misleading when $k>1$.  Fix $k\geq1$, at least $k+1$ participants, and $k+1$
distinct valid proposals $v_1,\ldots,v_{k+1}$.  Outputs are irrevocable partial
decision maps $m:P\rightharpoonup V$, ordered by extension, with at most $k$
values in their range.

At participant $p_r$'s local proposal input, the decision
$\{p_r\mapsto v_r\}$ is within spec.  The union of all proposal inputs is a common
input refinement.  Each local decision can continue there by itself, and any
family of at most $k$ distinct decisions can coexist.  The family of all $k+1$
cannot: every common refinement would preserve $k+1$ values and violate the
specification.  This is exactly the obstruction detected by Lifted
Monotonicity.  If a partition lets the owners decide independently, the
same family witnesses the negative side of Complete CAP.  Consensus is the
special case $k=1$.

This argument identifies the incompatibility in the specification.  Classical wait-free and
topological impossibility theorems additionally depend on a computation model,
failure assumptions, and progress requirements~\cite{saks2000set}.

\subsection{Membership-sensitive naming}

Tight strong renaming presents a different source of incompatibility.  The input
is the finite participant set $P$, ordered by inclusion.  The common output
universe contains finite partial naming maps; at input $P$, legal outputs are
injective maps
        \begin{equation}
  m:P\longrightarrow\{1,\ldots,|P|\},
  \label{eq:tight-renaming}
        \end{equation}
ordered by map extension.

Now start two processes independently.  At input $\{p\}$, tightness forces the
only name $p\mapsto1$; at input $\{q\}$ it likewise forces $q\mapsto1$.  The
combined input $\{p,q\}$ refines both singleton inputs.  Either assignment can
continue alone by giving the other process name 2, so Relational Monotonicity
is not the problem.  They cannot continue together: an injective map cannot
preserve both assignments to name 1.  Complete CALM fails on this two-row family,
and a deployment that places $p$ and $q$ in different components gives the
corresponding Complete CAP obstruction.

Notice why simply extending one completed naming is not a witness.  If $n$
processes already use names $1$ through $n$, a new process can receive $n+1$
without changing an old assignment.  The obstruction comes from independent
partial memberships making incompatible early choices.  The classical renaming
hierarchy adds model-specific topological and wait-free conclusions
~\cite{herlihy1991waitfree,herlihy1999topological}.

\subsection{What does a snapshot promise?}

The word ``snapshot'' covers interfaces with very different interface strength.
Consider inputs that are finite downward-closed causal histories, ordered by
induced extension.  A modest \emph{cut} interface outputs a finite
downward-closed subset of the available history as a lower bound on what has
happened.  Old cuts remain legal as history grows, and the union of finitely
many cuts remains downward closed.  This interface satisfies Lifted
Monotonicity and Directed Outputs.

That easy positive result should not be mistaken for a conventional snapshot
implementation.  The cut interface promises neither freshness nor completeness;
it could legally return an old cut.  A stronger interface may promise that a
vector or cut is the complete state at a globally finished boundary.  Under
equality of final answers, another relevant event changes the required result.
Functional CALM then says that the boundary must be sealed before the
completeness claim is exposed.

These are safety statements about two declared interfaces.  A distributed
snapshot protocol must also collect a consistent cut and terminate under its
channel and failure assumptions~\cite{chandy1985distributed}.  An atomic
snapshot records values at a chosen linearization point instead.  Classic
wait-free shared-memory implementations demonstrate that total ordering alone
does not prove a need for coordination; the process, memory, placement, and
failure model all matter~\cite{afek1993atomic}.

\begin{acks}
We thank Peter Alvaro, Natacha Crooks, Chris Douglas, Tyler Hou, Martin
Kleppmann, Paris Koutris, Edward Lee, Shulu Li, Vincent Liu, David Chu McElroy,
Sarah Morin, and Dan Suciu for helpful feedback on this work.

Generative AI tools (Claude and ChatGPT) were used as interactive writing and
reviewing assistants during the preparation of this work, including drafting
and revising text, checking mathematical arguments, and suggesting edits.  All
content was reviewed, validated, and approved by the author.
\end{acks}
